Unlearning audits plant canary records in a model's training data, request their deletion,
and test whether the deployed model still favours them.
Such an audit speaks for users' data only if the provider treats canaries like other
records. We study a provider that recognises canaries by their content and unlearns only them. For audits that average a bounded
per-pair statistic over randomly included canaries and excluded twins, we prove that the
leakage such a provider can hide, measured on records planted like canaries, is at most twice
the advantage of the best test that tells canaries from organic records. When leakage is
one-sided and determined by the trained record, the largest hideable leakage equals the power
of the optimal such test. A character n-gram classifier recognises both canary designs we
tested perfectly, including the random-string canaries of a recent anytime-valid
certificate. On Pythia-410M, a provider that unlearns only what the classifier flags obtains
the audit's certificate in 98\% of runs, while organic records with randomised inclusion keep
two thirds of their leakage, which honest unlearning removes. We prove that twins drawn from the organic record generator
keep the audit's exact null, and our detectors do not separate them beyond noise. At single
exposure, however, the audit then has no power and certifies a model that never unlearned. The bound shows that this trade-off is inherent: an audit can exceed organic leakage in
sensitivity only as far as its canaries are detectable. Membership-inference
audits with an unmatched holdout are invalid.
